\documentclass[11pt,a4paper]{article}
\usepackage{fontspec}
\usepackage[vietnamese]{babel}
\usepackage{geometry}
\geometry{a4paper, top=14mm, bottom=18mm, left=20mm, right=20mm, includehead, headheight=26pt, headsep=6mm}
\usepackage{amsmath,amssymb}
\usepackage{xcolor}
\usepackage{listings}
\usepackage{booktabs}
\usepackage{tcolorbox}
\usepackage{fancyhdr}
\usepackage{hyperref}

\pagestyle{fancy}
\fancyhf{}
\rhead{\textbf{Tuyển tập Lời giải Đề thi Olympic 11 TP.HCM}}
\lhead{Đội tuyển Olympic Tin học 11}
\rfoot{Trang \thepage}

\definecolor{codegreen}{rgb}{0,0.6,0}
\definecolor{codegray}{rgb}{0.5,0.5,0.5}
\definecolor{codepurple}{rgb}{0.58,0,0.82}
\definecolor{backcolour}{rgb}{0.96,0.96,0.96}

\lstdefinestyle{mystyle}{
    backgroundcolor=\color{backcolour},   
    commentstyle=\color{codegreen}\itshape,
    keywordstyle=\color{blue}\bfseries,
    numberstyle=\tiny\color{codegray},
    stringstyle=\color{codepurple},
    basicstyle=\ttfamily\footnotesize,
    breakatwhitespace=false,         
    breaklines=true,                 
    captionpos=b,                    
    keepspaces=true,                 
    numbers=left,                    
    numbersep=5pt,                  
    tabsize=4,
    frame=single
}
\lstset{style=mystyle}

\begin{document}

\begin{center}
    {\Large\bfseries TUYỂN TẬP LỜI GIẢI CHI TIẾT ĐỀ THI OLYMPIC TIN HỌC 11}\\[4pt]
    {\large\textbf{KHU VỰC THÀNH PHỐ HỒ CHÍ MINH (2024 -- 2026)}}\\[4pt]
    \textit{Phân tích thuật toán tối ưu \& Mã nguồn C++ chuẩn AC 100\% điểm số}
\end{center}

\vspace{0.4cm}

\tableofcontents

\newpage

\section{ĐỀ THI NĂM HỌC 2024 -- 2025}

\subsection{Bài 1: NGHE NHẠC (6.0 điểm)}
\textbf{File nộp:} \texttt{NGHENHAC.*} $\vert$ \textbf{Input:} \texttt{NGHENHAC.INP} $\vert$ \textbf{Output:} \texttt{NGHENHAC.OUT}

\paragraph{1. Đề bài tóm tắt}
An có $K$ phút rảnh rỗi và danh sách $N$ bài hát, bài thứ $i$ dài $m_i$ phút ($1 \le N \le 10^5, 1 \le K \le 10^9, 1 \le m_i \le 10^4$). Hãy tìm một \textbf{dãy liên tiếp} các bài hát sao cho số lượng bài hát là nhiều nhất và tổng thời lượng không vượt quá $K$.
\begin{itemize}
    \item Subtask 1 (60\% số điểm): $N \le 10^3$.
    \item Subtask 2 (40\% số điểm): $N \le 10^5$.
\end{itemize}

\paragraph{2. Phân tích thuật toán}
Sử dụng kỹ thuật \textbf{Hai con trỏ (Two Pointers)}:
\begin{itemize}
    \item Duy trì cửa sổ trượt $[left, right]$ và biến tổng thời gian \texttt{current\_sum}.
    \item Khi mở rộng $right$, nếu tổng thời gian vượt quá $K$, ta tăng $left$ để co ngắn cửa sổ.
    \item Độ phức tạp: $O(N)$ thời gian, $O(1)$ bộ nhớ phụ $\implies$ Đạt tối đa 6.0/6.0 điểm.
\end{itemize}

\begin{lstlisting}[language=C++]
#include <iostream>
#include <vector>
#include <algorithm>
using namespace std;

int main() {
    ios_base::sync_with_stdio(false);
    cin.tie(NULL);

    #ifndef ONLINE_JUDGE
    freopen("NGHENHAC.INP", "r", stdin);
    freopen("NGHENHAC.OUT", "w", stdout);
    #endif

    int n;
    long long k;
    if (!(cin >> n >> k)) return 0;

    vector<long long> m(n);
    for (int i = 0; i < n; i++) cin >> m[i];

    int left = 0, ans = 0;
    long long current_sum = 0;

    for (int right = 0; right < n; right++) {
        current_sum += m[right];
        while (current_sum > k && left <= right) {
            current_sum -= m[left];
            left++;
        }
        ans = max(ans, right - left + 1);
    }

    cout << ans << "\n";
    return 0;
}
\end{lstlisting}

\newpage

\subsection{Bài 2: HỌC TOÁN (7.0 điểm)}
\textbf{File nộp:} \texttt{HOCTOAN.*} $\vert$ \textbf{Input:} \texttt{HOCTOAN.INP} $\vert$ \textbf{Output:} \texttt{HOCTOAN.OUT}

\paragraph{1. Đề bài tóm tắt}
Cho hai số nguyên dương $a, b$ ($1 \le a < b \le 10^4$). Tìm hai số nguyên dương $x, y$ sao cho:
1. $\gcd(x, y) = a$
2. $\text{lcm}(x, y) = b$
3. Tổng $x + y$ nhỏ nhất. Dòng 1 ghi $x + y$, dòng 2 ghi $x, y$ ($x \le y$, chọn $x$ nhỏ nhất nếu trùng). Nếu không tìm được in $-1$.

\paragraph{2. Phân tích toán học \& Thuật toán}
Ta có đẳng thức cơ bản: $x \cdot y = \gcd(x, y) \cdot \text{lcm}(x, y) = a \cdot b$.
Đặt $x = a \cdot u, y = a \cdot v$ với $\gcd(u, v) = 1$. Khi đó:
\[
u \cdot v = \frac{b}{a}
\]
Điều kiện cần là $b$ phải chia hết cho $a$. Đặt $P = b / a$. Ta chỉ cần duyệt ước $u$ của $P$ từ $\lfloor\sqrt{P}\rfloor$ giảm dần về $1$. Cặp $(u, v)$ đầu tiên thỏa mãn $\gcd(u, v) = 1$ sẽ cho tổng $u + v$ nhỏ nhất. Độ phức tạp $O(\sqrt{b/a})$.

\begin{lstlisting}[language=C++]
#include <iostream>
#include <cmath>
using namespace std;

long long gcd_func(long long a, long long b) {
    while (b) { a %= b; swap(a, b); }
    return a;
}

int main() {
    ios_base::sync_with_stdio(false);
    cin.tie(NULL);

    #ifndef ONLINE_JUDGE
    freopen("HOCTOAN.INP", "r", stdin);
    freopen("HOCTOAN.OUT", "w", stdout);
    #endif

    long long a, b;
    if (!(cin >> a >> b)) return 0;

    if (b % a != 0) {
        cout << -1 << "\n";
        return 0;
    }

    long long P = b / a;
    long long best_u = -1, best_v = -1;

    for (long long u = sqrt(P); u >= 1; u--) {
        if (P % u == 0) {
            long long v = P / u;
            if (gcd_func(u, v) == 1) {
                best_u = u;
                best_v = v;
                break;
            }
        }
    }

    if (best_u == -1) cout << -1 << "\n";
    else {
        cout << a * (best_u + best_v) << "\n";
        cout << a * best_u << " " << a * best_v << "\n";
    }
    return 0;
}
\end{lstlisting}

\newpage

\subsection{Bài 3: ROBOT (7.0 điểm)}
\textbf{File nộp:} \texttt{ROBOT.*} $\vert$ \textbf{Input:} \texttt{ROBOT.INP} $\vert$ \textbf{Output:} \texttt{ROBOT.OUT}

\paragraph{1. Đề bài tóm tắt}
Bảng ô vuông kích thước $M \times N$ ($2 \le M, N \le 10^5$). Robot 1 đi từ $(1, 1)$ đến $(M, N)$ theo các lệnh \texttt{D} (xuống) và \texttt{R} (phải). Robot 2 đi từ $(M, N)$ về $(1, 1)$ theo các lệnh \texttt{U} (lên) và \texttt{L} (trái). Hãy đếm số ô mà cả hai robot cùng đi qua.

\paragraph{2. Phân tích thuật toán}
Kích thước bảng có thể lên đến $10^{10}$ ô nên tuyệt đối không được khai báo mảng 2 chiều. Mỗi robot chỉ đi tối đa $(M - 1) + (N - 1) \le 2 \cdot 10^5$ bước. Ta chỉ việc mô phỏng và lưu vết tọa độ vào \texttt{std::set<pair<int, int>>}. Độ phức tạp: $O((M + N) \log(M + N))$.

\begin{lstlisting}[language=C++]
#include <iostream>
#include <string>
#include <set>
using namespace std;

int main() {
    ios_base::sync_with_stdio(false);
    cin.tie(NULL);

    #ifndef ONLINE_JUDGE
    freopen("ROBOT.INP", "r", stdin);
    freopen("ROBOT.OUT", "w", stdout);
    #endif

    int m, n;
    if (!(cin >> m >> n)) return 0;
    string s1, s2;
    cin >> s1 >> s2;

    set<pair<int, int>> v1;
    int r = 1, c = 1;
    v1.insert({r, c});
    for (char ch : s1) {
        if (ch == 'D') r++;
        else if (ch == 'R') c++;
        v1.insert({r, c});
    }

    set<pair<int, int>> common;
    r = m; c = n;
    if (v1.count({r, c})) common.insert({r, c});
    for (char ch : s2) {
        if (ch == 'U') r--;
        else if (ch == 'L') c--;
        if (v1.count({r, c})) common.insert({r, c});
    }

    cout << common.size() << "\n";
    return 0;
}
\end{lstlisting}

\newpage

\section{ĐỀ THI NĂM HỌC 2025 -- 2026}

\subsection{Câu 1: DI TRUYỀN (6.0 điểm)}
\textbf{File nộp:} \texttt{DNA.*} $\vert$ \textbf{Input:} \texttt{DNA.INP} $\vert$ \textbf{Output:} \texttt{DNA.OUT}

\paragraph{1. Đề bài tóm tắt}
Cho $N$ chuỗi DNA ($1 \le N \le 10, |S| \le 10^5$). Một chuỗi là bất thường (đánh dấu 1, ngược lại 0) nếu thỏa ít nhất 1 trong 3 điều kiện:
\begin{enumerate}
    \item Có ít hơn 3 loại nucleotide khác nhau.
    \item Có ít nhất một loại chỉ xuất hiện đúng 1 lần.
    \item Khoảng cách giữa 2 nucleotide cùng loại liên tiếp lớn hơn ngưỡng $K$.
\end{enumerate}
Chuyển chuỗi nhị phân $N$ ký tự kết quả sang hệ thập phân và in ra.

\paragraph{2. Mã nguồn C++ hoàn chỉnh}
\begin{lstlisting}[language=C++]
#include <iostream>
#include <string>
#include <map>
using namespace std;

bool is_abnormal(const string &s, int k) {
    map<char, int> freq, last_pos;
    for (int i = 0; i < (int)s.size(); i++) {
        char c = s[i];
        freq[c]++;
        if (last_pos.count(c) && (i - last_pos[c]) > k) return true;
        last_pos[c] = i;
    }
    if (freq.size() < 3) return true;
    for (auto &p : freq) if (p.second == 1) return true;
    return false;
}

int main() {
    ios_base::sync_with_stdio(false);
    cin.tie(NULL);

    #ifndef ONLINE_JUDGE
    freopen("DNA.INP", "r", stdin);
    freopen("DNA.OUT", "w", stdout);
    #endif

    int n, k;
    if (!(cin >> n >> k)) return 0;
    long long ans = 0;
    for (int i = 0; i < n; i++) {
        string s; cin >> s;
        ans = ans * 2 + (is_abnormal(s, k) ? 1 : 0);
    }
    cout << ans << "\n";
    return 0;
}
\end{lstlisting}

\newpage

\subsection{Câu 2: THU HOẠCH (7.0 điểm)}
\textbf{File nộp:} \texttt{THUHOACH.*} $\vert$ \textbf{Input:} \texttt{THUHOACH.INP} $\vert$ \textbf{Output:} \texttt{THUHOACH.OUT}

\paragraph{1. Phân tích Quy hoạch động}
Cho $N$ thửa ruộng ($N \le 5000$). Cần chọn $K$ đoạn không giao nhau, mỗi đoạn độ dài đúng $L$ để tổng lớn nhất ($K \times L \le N$).
\begin{itemize}
    \item Mảng tiền tố tính tổng đoạn độ dài $L$: $val[i] = P[i] - P[i - L]$.
    \item Bảng phương án: $dp[i][j]$ là tổng lớn nhất khi xét đến thửa $i$ và chọn $j$ máy gặt:
    \[
    dp[i][j] = \max(dp[i - 1][j], dp[i - L][j - 1] + val[i])
    \]
\end{itemize}

\begin{lstlisting}[language=C++]
#include <iostream>
#include <vector>
#include <algorithm>
using namespace std;

const long long INF = 1e18;

int main() {
    ios_base::sync_with_stdio(false);
    cin.tie(NULL);

    #ifndef ONLINE_JUDGE
    freopen("THUHOACH.INP", "r", stdin);
    freopen("THUHOACH.OUT", "w", stdout);
    #endif

    int n, k, L;
    if (!(cin >> n >> k >> L)) return 0;

    vector<long long> pref(n + 1, 0), val(n + 1, 0);
    for (int i = 1; i <= n; i++) {
        long long w; cin >> w;
        pref[i] = pref[i - 1] + w;
    }
    for (int i = L; i <= n; i++) val[i] = pref[i] - pref[i - L];

    vector<vector<long long>> dp(n + 1, vector<long long>(k + 1, -INF));
    for (int i = 0; i <= n; i++) dp[i][0] = 0;

    for (int i = 1; i <= n; i++) {
        for (int j = 1; j <= k; j++) {
            dp[i][j] = dp[i - 1][j];
            if (i >= L && dp[i - L][j - 1] != -INF) {
                dp[i][j] = max(dp[i][j], dp[i - L][j - 1] + val[i]);
            }
        }
    }

    cout << dp[n][k] << "\n";
    return 0;
}
\end{lstlisting}

\newpage

\subsection{Câu 3: BỘ SƯU TẬP (7.0 điểm)}
\textbf{File nộp:} \texttt{SUUTAP.*} $\vert$ \textbf{Input:} \texttt{SUUTAP.INP} $\vert$ \textbf{Output:} \texttt{SUUTAP.OUT}

\paragraph{1. Đề bài tóm tắt}
Cho dãy $N$ viên đá ($N \le 10^5$), viên thứ $i$ có mã màu $C_i$ ($1 \le C_i \le 10^9$). Tìm độ dài đoạn liên tiếp ngắn nhất chứa ít nhất $K$ màu khác nhau.

\paragraph{2. Mã nguồn C++ Hai con trỏ + Bảng băm}
\begin{lstlisting}[language=C++]
#include <iostream>
#include <vector>
#include <map>
#include <algorithm>
using namespace std;

int main() {
    ios_base::sync_with_stdio(false);
    cin.tie(NULL);

    #ifndef ONLINE_JUDGE
    freopen("SUUTAP.INP", "r", stdin);
    freopen("SUUTAP.OUT", "w", stdout);
    #endif

    int n, k;
    if (!(cin >> n >> k)) return 0;
    vector<int> c(n);
    for (int i = 0; i < n; i++) cin >> c[i];

    map<int, int> freq;
    int left = 0, min_len = 1e9;

    for (int right = 0; right < n; right++) {
        freq[c[right]]++;
        while ((int)freq.size() >= k) {
            min_len = min(min_len, right - left + 1);
            freq[c[left]]--;
            if (freq[c[left]] == 0) freq.erase(c[left]);
            left++;
        }
    }

    if (min_len > n) cout << -1 << "\n";
    else cout << min_len << "\n";

    return 0;
}
\end{lstlisting}

\end{document}
